\addcontentsline{toc}{chapter}{Preamble}
\chapter*{Preamble}
\markboth{Preamble}{Preamble}

The complexity of systems that are relevant to engineering today has grown tremendously.
The control techniques based on frequency analysis that were perfectly adequate for simple systems tend to be difficult to use for more complex systems.
While the windup issues generated by ignoring constraints on the controller may be worked around in simple cases,
it is an indicator of the need to consider alternative techniques when state space and controller constraints play a fundamental role in a control problem.
On the other hand, techniques based on optimization such as Lyapunov functions and Model Predictive Control (MPC) may be considered overkill for simple systems
but they thrive in the complexity of modern systems as reported for instance in \cite{mayne2014model}:
\begin{quote}
  [The phenomenal success of MPC] in the process industries is well described in \cite{qin2003survey} and was mainly due to
  its conceptual simplicity and its ability to handle easily and effectively
  complex systems with hard control constraints and many inputs and outputs.
\end{quote}

An important challenge arising for these complex systems is the need to obtain sets satisfying given properties.
While the search over all sets is in general intractable, there is usually a well defined objective for the set to be computed that entirely depends on the application.
For instance, does the set need to outer or inner approximate a given set or does it need to satisfy a given property ?
If several sets are feasible, which one should be chosen ? The one with maximal/minimal volume, the one that contains a point the furthest possible in a given direction ?
How should the set be represented ? By inequalities that should be satisfied for the set ? Can it be represented as the projection of a simpler set ?

Once the intended end-use of the set as well as the properties it should satisfy are clarified,
a specific family of sets, often called \emph{template}, is chosen to formulate the search of the appropriate set as a problem that is numerically tractable.
The algorithms employed vary considerably depending on the choice of template but always largely rely on solving optimization problems.

In classical numerical optimization, the introduction of a modeling layer on top of optimization solvers has been a key enabler,
allowing different communities to share emergent technologies more easily.
As detailed in \cite{orchard-haysHistoryMathematicalProgramming1984}, it started in the early history of numerical optimization with a punch-card input format called MPS \cite{mps_ibm}.
Due to the lack of convenience of the MPS format,
algebraic modeling languages were developed such as AMPL \cite{Fourer1990} and GAMS \cite{brook1988gams}.
Later, high-level interpreted languages were developed that allowed to write numerical
algorithms in a high-level language, while the core of the methods used were implemented
in low-level compiled languages such as C or Fortran, e.g. the optimized LAPACK~\cite{anderson1999lapack} library for numerical computing.
For numerical algorithms written in these high-level languages that needed to
solve optimization programs, it became essential to be able to access algebraic modeling frameworks embedded in the language.
This led to the creation CVX~\cite{grant2014cvx} and YALMIP~\cite{lofberg2004yalmip} in MATLAB\texttrademark{} and CVXPY~\cite{diamond2016cvxpy} and Pyomo~\cite{hart2017pyomo} in Python.
For large-scale problems and problems requiring substantial transformations,
the interpreted nature of the high-level languages led to significant time spent in model transformation compared to the solve time.
This issue was resolved by the creation of the high-level compiled language Julia
and the \jlpkg{Convex}~\cite{convexjl} and \jlpkg{JuMP}~\cite{dunning2017jump} algebraic
modeling languages embedded in Julia that are built on top of
the MathOptInterface~\cite{legat2020mathoptinterfacearXiv} layer for interfacing with the solvers.

When the purpose is to compute sets rather than numbers, the numerical complexity is amplified,
yet there does not currently exist such a clearly defined interface.
The purpose of this thesis is to define this interface between
\begin{enumerate}
  \item algorithms relying on the computation of the solution of a \emph{set program}, that is, the computation of a set satisfying given properties and maximizing a given objective, these are independent of the way these sets are computed; and
  \item algorithms for solving a specific class of set program for a given template,
    these are used as interchangeable solvers for solving the set programs
    formulated by algorithms described the previous item.
\end{enumerate}

The set program searching over all sets, not specific to any template, is called the \emph{generic set program}.
The generic set program restricted to sets of a given template $T$ is called the \emph{set program for $T$}.
Several templates are studied in this thesis including the classical templates polyhedra, zonotopes and ellipsoids
but also more elaborate ones that provide a richer family of sets but are more complicated to implement both theoretically and algorithmically.
These includes sublevel sets of polynomial forms, sets with support function that is the root of a polynomial,
piecewise semi-ellipsoidal sets, and piecewise polynomial sets.

We study two questions in detail about set programs that both examine different aspects of duality.

\begin{itemize}

\item \textit{Conic duality} \\

First, we reinterpret the classical duality of conic programs for set programs.
That is, given a definition of the convex combination of two sets, a template can be regarded as a convex cone where each point of the cone is a set of this template.
As the set program for a given template is a restriction of the generic set program,
its dual is a relaxation of the dual of the generic set program.
This means for instance that while an infeasibility certificate of the dual of a generic set program
is a valid certificate for any set,
an infeasible certificate of the dual of a set program for a given template is only
valid for this template.
In other words, it is a relaxed infeasibility certificate for the generic set program.

In numerous fields, the solution of relaxations can still be ``rounded'' to an actual solution.
This is for instance the purpose of \emph{heuristic callbacks} in mixed integer programming,
the \emph{randomized hyperplane rounding} technique for Max-Cut~\cite{goemans1995improved}
or the \emph{low rank rounding} of moment matrices to moments of Dirac measures in polynomial optimization; see \secref{atom} for more details.
While the rounding techniques can be shown to work reliably on specific instances,
a given guarantee on the quality of the rounded solution is appreciated
as it ensures that an algorithm can rely on the performance of the rounding.
We provide in \theoref{fourth} a rounding of the infeasibility certificate
of a certain class of set program for the family of sublevel sets of polynomial forms
as well as a guarantee on the quality of the rounded infeasibility certificate for the generic set program.

\item \textit{Geometric duality} \\

Second, we focus on decision variables of set program that are constrained to be convex sets.
In this case, the set can be represented in two different ways that take various forms depending on the template.
\begin{enumerate}
  \item The first way is with the \emph{gauge} or \emph{Minkowski} function of the set.
    In other words, the set is represented by several inequalities that all points in the set should satisfy.
    For polyhedra, this is called the H-representation.
    For ellipsoids, this is the positive definite matrix $Q$ such that the ellipsoid is the set of points $x$ such that $x^\Tr Q x \le 1$.
  \item The second way is with the \emph{support} function of the set.
    This represents the \emph{polar} set by several inequalities that all points in the polar set should satisfy.
    For polyhedra, this is called the V-representation.
    It represents the set as the convex hull of points.
    For ellipsoids, this is the positive definite matrix $Q$ such that the ellipsoid is the set of points $x$ such that $x^\Tr Q^{-1} x \le 1$
    and the polar of the ellipsoids is the set of points $y$ such that $y^\Tr Q y \le 1$.
\end{enumerate}
An algorithm computing a convex set of a given template needs to choose the representation that is used for computation.
Indeed, changing the representation is often not an option.
For instance, computing the V-representation from the H-representation of a polyhedron is computationally expensive
and when formulating a semidefinite program that searches for an ellipsoid $\mathcal{E}_Q$, the decision variable is either the matrix $Q$ or the matrix $Q^{-1}$.
Moreover, the choice of representation for a given set may affect the choice for another set.
For instance, given an inclusion constraint between two sets, if both sets are represented with their support function, this can be formulated
as an inequality between the support function but if one set is represented with its support function and the other one with its gauge function,
it is not clear how the inclusion can be encoded in terms of the support and gauge functions\footnote{except in specific cases, for instance checking the inclusion of a polyhedron given its V-representation in a polyhedron given its H-representation simply amounts to checking whether each element of the V-representation of the first set is inside all halfspaces in the H-representation of the second set, see \propref{vinhrep}.}.

The key question is therefore: which representation should be used for each set in the set program ?
We address this question for specific cases of inclusion constraints in \theoref{gaugesupp_image} and \theoref{controlineq} and highlight the fact that this answer is essentially independent from the template.
The choice is fundamentally based on convex properties of gauge or support functions and on the type of inclusion constraint the sets are involved in.
This sheds a light on this choice of representation which often resort to template-specific arguments.


\end{itemize}





The thesis is organized as follows.
The first part introduces background materials while the second part builds on these results to formulate our contributions.
The first part is divided into three chapters.
\chapref{conv} provides the background on convex analysis as well as algebraic geometry needed for the thesis.
Then, \chapref{opti} surveys optimization results and algorithms that the thesis relies on.
Finally, \chapref{sys} introduces the classes of systems under consideration as well as classical techniques used to compute invariant sets for these systems.

The second part is split into four chapters.
In \chapref{setprog}, we define the set programming interface.
We first discuss the different formulations of generic set programs
as well as the choice of representation for convex sets.
Then we particularize these results to the different templates considered in the thesis.
In \chapref{stab}, we explore the conic duality of set programs and provide a rounding procedure to deduce infeasibility
certificates for the generic set program given an infeasibility certificate for the set program for sublevel sets of polynomials.
This can be used to certify the unstability of a switched system.
We provide as well a guarantee of the quality of the certificates that is based on the entropy of the language generated by the automaton formed by the constraints of the set program.
In \chapref{mci}, we use the geometric duality of set programs to compute controlled invariant sets for control systems.
This is applied to model predictive control and stochastic programming.
In \chapref{entropic}, we show how to formulate an approximation
of the entropic cone as a hierarchy of set programs.
We then show how to circumvent the curse of dimensionality in the high-dimensional
polyhedral approximation of the entropic cone by approximating it
only in a specific direction using dual dynamic programming.

\paragraph{Reproducibility}
We put a particular emphasis on the reproducibility of our result
and the quality and reusability of the code.
The codes used to obtain the results of \secref{entropyguarantee}, \secref{certifylowerbound} and \exemref{simplesetprog}
are published on codeocean respectively in \cite{c2816dff-2bf4-43b8-bb71-f78ab352d9ae}, \cite{legat2020certifyingCO} and \cite{legat2020piecewiseCO}.

These codes rely on both new Julia packages we have developed and existing Julia packages.
We now briefly mention the packages developed for this thesis in collaboration with several other contributors that we mention in the acknowledgment. %
Julia is a programming language particularly well suited for numerical computing.
Its \emph{multiple dispatch} mechanism as well as its \emph{just-in-time} compilation nature
make it particularly well suited for numerical computing~\cite{bezanson2017julia}.

Related to \secref{polyhedra}, the package \jlpkg{Polyhedra} provides a unified interface for polyhedral computation library.
Currently, CDD~\cite{fukuda2003cdd}, LRS~\cite{avis2000revised} and QHull~\cite{barber1996quickhull} are supported through the \jlpkg{CDDLib}, \jlpkg{LRSLib} and \jlpkg{QHull} packages respectively. The \jlpkg{ConvexHull} pure julia library also implements the \jlpkg{Polyhedra} interface.
The main features of \jlpkg{Polyhedra} are representation conversion, polyhedral projection, intersections, convex/conic hulls of union, Minkowski sum, linear image and preimage, redundancy removal and 2D visualization with \jlpkg{Plots} and 3D visualization using \jlpkg{Makie} or \jlpkg{MeshCat}.

Related to \secref{algebraic}, the package \jlpkg{MultivariatePolynomials} provides a unified interface for implementation of polynomials.
Two implementations are available: \jlpkg{DynamicPolynomials} and \jlpkg{TypedPolynomials}.
The package \jlpkg{SemialgebraicSets} implements datastructures for basic semialgebraic sets
and varieties and implement the Buchberger's algorithm
as well as the approach presented in~\cite{manocha1995algorithms} for computing the
elements of a zero-dimensional variety from the multiplication matrices; see \secref{solvealg}.

Related to \chapref{opti}, \jlpkg{MathOptInterface} defines an interface
for communicating with optimization solvers implementing the
bridges mechanism described in \secref{bridges}
and \jlpkg{JuMP} provides an algebraic modeling language on top of this interface.
We developed the MathOptInterface for the following SDP solvers:
\jlpkg{CDCS}~\cite{zheng2020chordal},
\jlpkg{CSDP}~\cite{borchers1999csdp},
\jlpkg{DSDP}~\cite{dsdp1},
\jlpkg{Mosek}~\cite{mosek2017mosek81082},
\jlpkg{SCS}~\cite{ocpb:16},
\jlpkg{SDPA}~\cite{yamashita2003implementation},
\jlpkg{SDPT3}~\cite{tutuncu2003solving},
\jlpkg{SDPNAL}~\cite{yang2015sdpnal},
and \jlpkg{SeDuMi}~\cite{sturm1999using}

Related to \secref{sosprog}, we developed
\jlpkg{MultivariateBases} that provides an interface for bases of multivariate polynomials as well as the implementation of some of them,
\jlpkg{MultivariateMoments} that implements datastructures for moment series and moment matrices as well as the atom extraction procedure described in \secref{atom} based on the method developed in \jlpkg{SemialgebraicSets} for finding the elements of zero-dimensional varieties.
\jlpkg{PolyJuMP} that extends \jlpkg{JuMP} and \jlpkg{MathOptInterface} to polynomial variables, polynomial equalities and polynomial inequalities.
\jlpkg{SumOfSquares} that extends \jlpkg{JuMP} and \jlpkg{MathOptInterface} to sum-of-squares variables, sum-of-squares constraints.

Related to \secref{ddp}, we developed \jlpkg{StructDualDynProg} that implements dual dynamic programming on arbitrary markov chains with support for both optimality and feasibility cuts. If too many cuts start getting accumulated, the package can eliminate cuts based on different heuristics. These heuristics are implements in \jlpkg{CutPruners}.

Related to \chapref{sys}, we developed \jlpkg{MathematicalSystems} that defines an interface
for dynamical systems and implements datastructures for the commonl used ones.
Related to \secref{hybrid}, we developed \jlpkg{HybridSystems} that extends \jlpkg{MathematicalSystems} to hybrid systems.

Related to \chapref{setprog}, we implemented \jlpkg{SetProg} that extends \jlpkg{JuMP}
to set variables and constraints and implements the reformulation to \jlpkg{SumOfSquares}
in two steps: first determine the representation of the set variables, second reformulate the objective and each constraint independently.
The template of each variable is specified at the creation of the variable and the rest
of the set program is specified independently of the template.

Related to \chapref{stab} and \chapref{mci}, we implemented \jlpkg{SwitchOnSafety}
that compute invariant sets, from a given template, of systems defined with \jlpkg{MathematicalSystems} or hybrid systems defined with \jlpkg{HybridSystems}
using \jlpkg{SetProg}.
It also implements the alternative techniques for approximating the joint spectral radius
described in \secref{certifylowerbound}.

Related to \chapref{entropic},
we developed \jlpkg{EntropicCone} for working with the \tEc{}.
It implements optimization on the entropic cone, as discussed in \secref{optientropic},
using the implementation of \algoref{dirobj} included in \jlpkg{StructDualDynProg}.

The codes of \exemref{4normpolar}, \exemref{newton}, \exemref{moment} and \exemref{atom}
using the above mentioned packages are included inline in the text for ease of reproducibility.

\newpage

\section*{Bibliographic notes}

\subsection*{Conference Publications}

\begin{enumerate}

\item \myfullcite{legat2016generating}
\item \myfullcite{gomes2017stable}
\item \myfullcite{legat2018computing}
\item \myfullcite{gomes2018minimally}

\end{enumerate}


\subsection*{Journal Publications}

\begin{enumerate}
\item \myfullcite{legat2019entropy}
\item \myfullcite{legat2020sum}
\item \myfullcite{legat2020certifying}
\item \myfullcite{legat2020piecewiseCO}
\end{enumerate}
